\documentclass[12pt,a4paper]{report}
\usepackage[utf8]{inputenc}
\usepackage[T1]{fontenc}
\usepackage[french]{babel}
\usepackage{graphicx}
\usepackage{hyperref}
\usepackage{booktabs}
\usepackage{geometry}
\geometry{margin=2.5cm}

\title{Projet de Traducteur Multilingue Français--Anglais--Éwé avec Reconnaissance Vocale}
\author{%
  \textbf{Groupe GR02} \\[0.4em]
  AZIGLOSSOU Judith Gracia \\
  NAPO NISSAO Rafia \\
  WOGLO Koffi Boris Gilbride \\
  PIO Sikirou Afeez Prince Cyrille \\
  KOMBATE Gountande \\[0.6em]
  \small Master 1 IA-BD EPL S2 --- Informatique et Gestion de données
}
\date{Juin 2026}

\begin{document}
\maketitle

\begin{center}
\fbox{\parbox{0.85\textwidth}{\centering
\textbf{Travail de groupe --- GR02} \\
Membres : AZIGLOSSOU Judith Gracia, NAPO NISSAO Rafia, WOGLO Koffi Boris Gilbride,
PIO Sikirou Afeez Prince Cyrille, KOMBATE Gountande.}}
\end{center}

\tableofcontents

\chapter*{Résumé}
\addcontentsline{toc}{chapter}{Résumé}
Ce projet présente la conception d'un système de traduction automatique multilingue couvrant le français, l'anglais et l'ewe, avec une extension vocale. L'approche proposée s'appuie sur les modèles Transformer multilingues, en particulier NLLB-200, et tient compte des contraintes propres aux langues africaines à faibles ressources.

\chapter{Introduction}
La traduction automatique constitue un enjeu majeur d'inclusion numérique. L'ewe, langue parlée au Togo et au Ghana, reste peu représentée dans les grands corpus numériques. Ce travail propose une architecture réaliste associant collecte de données, nettoyage NLP, fine-tuning, API et interface utilisateur.

\chapter{Méthodologie}
La méthodologie comprend la collecte de corpus parallèles, la normalisation Unicode, le fine-tuning d'un modèle NLLB-200 distilled 600M avec LoRA, et l'évaluation par BLEU, chrF, WER et CER.

\chapter{Architecture du système}
Le système est constitué de quatre couches :
\begin{itemize}
  \item une \textbf{interface React/Vite} (saisie texte, traduction instantanée, historique et favoris, dictée et lecture vocale) ;
  \item une \textbf{API FastAPI} exposant \texttt{/translate}, \texttt{/transcribe}, \texttt{/synthesize}, \texttt{/health} et \texttt{/about} ;
  \item un \textbf{moteur de traduction} NLLB-200-distilled-600M fine-tuné avec LoRA, avec repli automatique sur un dictionnaire de démonstration en cas d'indisponibilité du modèle ;
  \item un \textbf{module vocal} : reconnaissance vocale via l'API Web Speech du navigateur (français, anglais) et Whisper côté serveur (activable), synthèse vocale du navigateur pour la traduction orale.
\end{itemize}

\chapter{Résultats et discussion}
Le fine-tuning LoRA (rang 16, $\alpha=32$, projections \texttt{q\_proj}/\texttt{v\_proj}) de NLLB-200 a été évalué par BLEU et chrF (sacreBLEU) sur un jeu de test multidirectionnel.

\begin{table}[h]
\centering
\begin{tabular}{llrr}
\toprule
Source & Cible & BLEU & chrF \\
\midrule
Anglais & Français & 38.18 & 64.21 \\
Français & Anglais & 42.75 & 62.02 \\
Éwé & Français & 17.20 & 37.38 \\
Français & Éwé & 9.33 & 30.32 \\
\bottomrule
\end{tabular}
\caption{Scores BLEU/chrF par direction de traduction.}
\end{table}

Les paires français--anglais atteignent une qualité élevée, conforme aux capacités multilingues de NLLB. Les directions impliquant l'éwé restent plus faibles, ce qui est attendu pour une langue à faibles ressources : les scores chrF (37 et 30) montrent néanmoins une couverture lexicale réelle. Les limites majeures concernent la quantité et la qualité des corpus alignés éwé (forte présence de données religieuses, à équilibrer) et la durée d'entraînement.

\chapter{Conclusion}
Le projet constitue une base crédible pour explorer la traduction automatique des langues africaines à faibles ressources dans un contexte universitaire.

\bibliographystyle{plain}
\begin{thebibliography}{9}
\bibitem{nllb} Meta AI. No Language Left Behind: Scaling Human-Centered Machine Translation. 2022.
\bibitem{lora} Hu et al. LoRA: Low-Rank Adaptation of Large Language Models. 2021.
\bibitem{sennrich} Sennrich et al. Improving Neural Machine Translation Models with Monolingual Data. 2016.
\end{thebibliography}
\end{document}
